\section{Conclusion}
\label{sec:conclusion}

The curse of dimensionality is not a law of nature. It is a consequence of spherical geometry---specifically, the volume concentration of the hypersphere in high dimensions. By operating on a different manifold, the curse can be broken.

We have shown that the Involuted Oblate Toroid (IOT) metric space reverses the dimensional volume collapse through two mechanisms: an exponential volume gradient between the torus's outer and inner edges, and an involution symmetry that decorrelates pairwise distances by introducing a per-dimension binary path choice. These mechanisms combine to produce a critical scaling law:
\begin{equation*}
    \left(\frac{r}{R}\right)_{\text{critical}} = \frac{\ln d}{2d}
\end{equation*}
that cleanly separates dimensional contraction (the curse) from dimensional expansion (the anti-curse).

The result traces a line from Archimedes' sphere-cylinder correspondence through the topology of the torus to a constructive resolution of a problem that has shaped decades of algorithm design. The sphere's volume can be understood via projection onto a cylinder. The cylinder, with identified ends, becomes a torus. The torus, with coupled radii and an involution, becomes a manifold where the curse runs in reverse.

At $R{:}r = 30{:}1$ and $d = 308$ dimensions, the IOT metric maintains a distance contrast ratio of 13.7:1---meaningful nearest-neighbor discrimination at a dimensionality where Euclidean distance has fully degenerated. The dimension-adaptive formulation $r/R = c \cdot \ln(d)/d$ extends this property to arbitrary dimensionality.

The curse was never about dimensionality. It was about geometry. Change the geometry, and the curse becomes a blessing.
